\section{Operadores fundamentales del TPK: dominios, acciones, invariantes y
memoria}
\label{sec:inventario-operatorio-tpk-rev11}
\index{TPK!inventario operatorio}

La \cref{def:categorias-operatorias-tpk-rev11} ha fijado primero las ocho
clases ontológico--operatorias: soporte y estado; selección; transporte;
lectores y cocientes; memoria y frontera; periodicidad y retorno;
prolongación; y salidas. Esta sección define sus realizaciones. Para cada
operador se declara la séxtupla \((D,C,f,I,L,M)\) de dominio, codominio, regla,
invariantes, información publicada o perdida y memoria conservada. El censo
nominal posterior agrupa símbolos ya definidos por esos seis datos y no
sustituye sus dominios ni sus reglas.

\begin{definition}[Ficha tipada de un operador del TPK]
Un operador canónico es una séxtupla
\[
 \mathcal O=(D_{\mathcal O},C_{\mathcal O},f_{\mathcal O},
 I_{\mathcal O},L_{\mathcal O},M_{\mathcal O}),
\]
donde \(f_{\mathcal O}:D_{\mathcal O}\to C_{\mathcal O}\) es la regla de
actuación, \(I_{\mathcal O}\) es la familia de relaciones que conserva,
\(L_{\mathcal O}\) declara las coordenadas publicadas y la relación de
equivalencia inducida por el olvido, y \(M_{\mathcal O}\) enumera los datos
que permanecen en el estado enriquecido después de la publicación. Dos
realizaciones representan el mismo operador únicamente cuando existe una
equivalencia tipada que entrelaza sus reglas y sus invariantes; compartir
periodo, cardinal o salida visible no basta.
\end{definition}

\subsection{Lectores modulares construidos}

La reducción nonádica se fija por el representante publicado
\begin{equation}
 \rho_9(n)=1+((n-1)\bmod 9)
 =n-9\left\lfloor\frac{n-1}{9}\right\rfloor,
 \label{eq:rho9-construida-tpk-rev11}
\end{equation}
de modo que la clase nula se escribe como \(9\). Para \(n\geq1\), la raíz
digital satisface
\[
 \operatorname{dr}_9(n)=\rho_9(n),
\]
pero las dos expresiones conservan semánticas distintas: la primera es una
sección del cociente; la segunda publica la suma iterada de cifras. El levantamiento
trítico escribe de manera única
\begin{equation}
 n=r_3(n)+3c_3(n),
 \qquad r_3(n)\in\{-1,0,+1\},\quad c_3(n)\in\mathbb Z,
 \label{eq:lift-mod3-tpk-rev11}
\end{equation}
y conserva simultáneamente residuo orientado y cociente de acarreo. Por
último, el lector de bloques decimales es la división euclídea
\begin{equation}
 N=1000q_{1000}(N)+r_{1000}(N),
 \qquad 0\leq r_{1000}(N)<1000,
 \label{eq:lector-mod1000-tpk-rev11}
\end{equation}
por lo que base \(1000\), residuo módulo \(1000\) y acarreo \(q_{1000}\)
son tres datos relacionados y no un mismo objeto.

\subsection{Clasificación de las realizaciones operatorias por dominio y función}

Las definiciones anteriores gobiernan la lectura de la tabla. Sus catorce
agrupaciones funcionales refinan las ocho clases operatorias y resumen las
realizaciones que comparten función; una clase puede, por tanto, ocupar más
de una fila. La tabla no opera como definición sustitutiva. Los símbolos que
comparten un cardinal permanecen separados por dominio y codominio.

\begingroup
\footnotesize
\renewcommand{\arraystretch}{1.06}
\begin{longtable}{@{}p{.21\textwidth}p{.31\textwidth}p{.40\textwidth}@{}}
\toprule
Agrupación funcional & Operadores y espacios & Dominio, acción e invariantes \\
\midrule
\endfirsthead
\toprule
Agrupación funcional & Operadores y espacios & Dominio, acción e invariantes \\
\midrule
\endhead
Estado y memoria de ruta &
\(\mathcal X_{\rm TPK}^{\rm enr}\), fibra \(\mathcal F_q\), memoria de ruta &
Conservan posición, hoja, orientación, residuo, cociente, acarreo, firma,
frontera, cilindro, prefijo, genealogía y memoria. \\

Selección de hoja &
\(M_{\rm ph}:\{1,\ldots,9\}\to\{-1,0,+1\}\) &
Activa la evaluación aditiva, la multiplicativa o el paso de umbral; no
mueve por sí mismo el estado. \\

Transporte espacial &
\(T_N,T_E,T_S,T_O\), motor reversible de dos cursores,
\(F_{\rm obs}\) &
Realiza las traslaciones del grafo de Cayley, actualiza hoja y orientación y
conserva cada transición de la ruta. \\

Lectores modulares &
\(\rho_9\), \((r_3,c_3)\), \((q_{1000},r_{1000})\) &
Publican fase nonádica, orientación trítica y bloque decimal sin borrar los
cocientes de levantamiento. \\

Bloque y memoria &
\(\mathcal K_{27}=F_{\rm obs}^{27}\), filtro de memoria
\(\mathcal N_{27}\) &
El primero compone veintisiete transportes; el segundo selecciona los
eventos de memoria. Igual cardinal no implica igual operador. \\

Estado dodecafásico &
\(\mathcal R_{12}\), centros \(\theta_m=30m-10\) grados,
refinamiento \(\mathcal R_{120}\) &
Organiza las doce fases, su oposición y el refinamiento angular en incrementos de
tres grados. \\

Periodicidad y profundidades &
\(C_{12}\hookrightarrow C_{108}\twoheadrightarrow C_9\),
\(\Pi_{108}^{\rm cur}\), \(w_{108}\), \(w_{1080}\),
\(\Gamma_0(1080)\) &
Distingue la periodicidad total de orden \(108\), la proyección de cursores, las
palabras de 108 y 1080 trits y el nivel modular 1080. \\

Canales (90/120) &
máscara \(C_m\), \(S_{90},S_{120}\), gramática orbital,
supertraza e invariantes \(\nu_{120},\nu_{270}\) &
Separa las dos genealogías, sus colas geométricas, sus pesos y las
coordenadas que extraen del estado. \\

Bidireccionalidad &
\(P_\pm\), inversión de ruta, \(\operatorname{Ext}_\pm\),
\(N_\pm\), \(\beta\), \(C_M=-\operatorname{rev}\) &
Realiza, respectivamente, intercambio espectral, ida y retorno espacial,
prolongación futura/pasada, valencia, balance y conjugación de palabra. \\

Emisión y coinducción &
\(L_0,L_1\), \(R_{36}\), \(G_9\), partición \(3^6=729\),
\(\varprojlim\operatorname{Cyl}_m\) &
Eleva la semilla, abre la frontera y enumera exhaustivamente las
prolongaciones finitas. El límite numérico y el levantamiento enriquecido
pertenecen a tipos distintos. \\

Lectores numéricos &
\(\mathscr C_\pi\), \(\mathscr P_e\), \(F_{\rm av}\) &
Construyen, por horquillas racionales, recurrencia factorial y autoescala de
Fibonacci, las secciones publicadas de \(\pi,e,\varphi\). \\

Proyección excepcional &
levantamiento de Witt, enumerador de peso y sombra no selectora &
Transporta el mismo estado a incidencia ternaria sin intervenir en la
selección arquimediana de cifras. \\

Atlas y realización &
firma de ruta, atlas (81/56/13/324), involución \(C_M\), centro
electrónico y operadores de fibra &
Convierte el estado de ruta en clasificadores, caracteres y operadores antes de
aplicar cualquier carta dimensional. \\

Validación finita &
censos exhaustivos, ablaciones e independencia causal &
Comprueba cobertura y unicidad en la profundidad calculada, y verifica que
ningún prefijo objetivo interviene en el mapa generador. \\
\bottomrule
\end{longtable}
\endgroup

\Needspace{46\baselineskip}
\subsection{Índice funcional de los operadores canónicos}

La clasificación anterior organiza las realizaciones por función. La tabla siguiente
reúne los operadores después de la definición categorial y de los lectores
construidos. Cada fila declara dominio, codominio, regla e invariantes; los
símbolos permiten reconocer una misma definición a lo largo de sus
distintas realizaciones sin identificar operadores de tipos diferentes.

\begingroup
\footnotesize
\renewcommand{\arraystretch}{1.06}
\enlargethispage{2\baselineskip}
\begin{longtable}{@{}r >{\raggedright\arraybackslash}p{.42\textwidth}
  >{\raggedright\arraybackslash}p{.22\textwidth}
  >{\raggedright\arraybackslash}p{.22\textwidth}@{}}
\toprule
\# & Operador canónico & Símbolo & Capa tipada \\
\midrule
\endfirsthead
\toprule
\# & Operador canónico & Símbolo & Capa tipada \\
\midrule
\endhead
1 & Célula aritmética aditivo--multiplicativa de orden nueve
  & \(\mathcal A_9\) & soporte aritmético \\
2 & Descomposición residual y cociente nonádicos
  & \((\rho_9,q_9)\) & reducción y elevación \\
3 & Reducción ternaria orientada
  & \(\rho_3^{\rm or}\) & reducción y elevación \\
4 & Lector posicional ternario--decimal de base mil
  & \(\iota_{1000}\) & reducción y elevación \\
5 & Acción cardinal orientada sobre el toro aritmético
  & \(T_d,\ d\in\{N,E,S,O\}\) & transporte cardinal \\
6 & Realimentación local aditivo--multiplicativa
  & \(F_{\rm loc}\) & transporte cardinal \\
7 & Emisor mínimo de dos cursores
  & \(T_{\min}\) & emisión finita \\
8 & Acoplador decimal de firmas aritméticas
  & \(G\) & emisión finita \\
9 & Proyección ternaria del catálogo
  & \(\Psi_6\) & emisión finita \\
10 & Sobreyectividad local del transporte cardinal
  & \(\mathcal R_3\) & transporte cardinal \\
11 & Lazo de transporte de veintisiete iteraciones
  & \(\mathcal K_{27}\) & transporte cardinal \\
12 & Estado enriquecido del núcleo topológico de fase
  & \(\mathcal X_{\rm TPK}^{\rm enr}\) & estado enriquecido \\
13 & Transición de estado Hensel
  & \(\mathcal H\) & estado enriquecido \\
14 & Transición exacta de cilindros ternario--decimales
  & \(\mathcal C_{3,1000}\) & estado enriquecido \\
15 & Firma conjunta de frontera
  & \(\Sigma_{\rm B}\) & estado enriquecido \\
16 & Relación de extensión y supervivencia nonádica
  & \({\rm EF}\text{--}G_9\) & prolongación y monodromía \\
17 & Monodromía nonádica con holonomía de frontera
  & \(\mathcal M_9\) & prolongación y monodromía \\
18 & Sección global del sistema inverso de supervivientes
  & \(X_\chi=\varprojlim_m{\rm Surv}^{(\chi)}_m\)
  & prolongación y monodromía \\
19 & Categoría bivariada de caminos aritméticos
  & \(\mathcal P_{\rm APP}^{(2)}\) & transporte cardinal \\
20 & Cronología observable de dos cursores
  & \(F_{\rm obs}\) & transporte cardinal \\
21 & Alfabeto decimal de cuarenta y dos tríadas
  & \(\mathcal A_{42}\) & emisión finita \\
22 & Relación tipada de extensión de fibras
  & \({\rm Ext}_r\) & estado enriquecido \\
23 & Transferencia nonádica de fibras compatibles
  & \(K_{\rm ph}\) & prolongación y monodromía \\
24 & Refinamiento de semillas por avance y giro
  & \((P,H,Q)\) & transporte cardinal \\
25 & Cociclos de modo y orientación
  & \((c_{\rm mode},c_{\rm or})\) & estado enriquecido \\
26 & Prolongación triádica de veinte bloques
  & \(E_{20}\) & prolongación y monodromía \\
27 & Sistema dodecafásico orientado
  & \(\mathcal R_{12}\) & lectura dodecafásica \\
28 & Observable armónico dodecafásico firmado
  & \(U_{12}^{\rm sgn}\) & lectura dodecafásica \\
29 & Transformación armónica por órbitas de paso tres
  & \(\mathcal H_{12}\) & lectura dodecafásica \\
30 & Extractor transversal de diferencias cíclicas
  & \(\mathcal E_{3,4}\) & lectura dodecafásica \\
31 & Recurrencia dodecafásica de cierre en base mil
  & \(\mathcal C_\alpha\) & cierre en base mil \\
32 & Transductor de historia enriquecida a carácter angular
  & \((L_{\rm tick},\chi_{\rm mass})\) & interfaz dimensional \\
33 & Retícula bidireccional de índice doce
  & \(\mathcal W_{\pm}\) & interfaz dimensional \\*
34 & Atlas nonádico de familias de extensión
  & \(\mathcal A_{\rm species}\) & interfaz dimensional \\
\bottomrule
\end{longtable}
\endgroup

El símbolo \(L_{\rm tick}\) designa el lector de iteración y retorno. Su
codominio es la memoria temporal del estado; no define una unidad temporal
física independiente.

La función selectora interna
\[
 M_{\rm ph}:\mathcal D_9\longrightarrow\{+1,-1,0\}
\]
pertenece al TPK\@; su vector tabulado es
\[
 m_{\rm ph}:=\bigl(M_{\rm ph}(1),\ldots,M_{\rm ph}(9)\bigr)^{\mathsf T}.
\]
La función \(M_{\rm ph}\), el vector \(m_{\rm ph}\) y el operador de estado
son objetos de tipos distintos. El valor escalar \(M_{\rm ph}(g(t))\) entra
como segundo argumento del selector \(\operatorname{Sel}_t\); no se compone
directamente con el estado. El transporte cardinal \(T_d\) ocupa otro factor
de la transición \(\mathcal U_t\) y realiza el movimiento sobre APP mediante
su elevación \(\operatorname{Lift}(T_d)\). Esta separación conserva un único
contenedor, distingue selección de fase y transporte, y preserva la
arquitectura completa por tipos.

\Needspace{9\baselineskip}
\begin{resultbox}[title={Profundidad finita y límite enriquecido}]
La enumeración, los censos y la contracción de cilindros numéricos son
resultados exactos en sus dominios. La existencia y unicidad de una sección
enriquecida sometida conjuntamente a todos los predicados de supervivencia
se obtiene cuando las fibras de levantamiento son no vacías, unitarias y
compatibles en toda profundidad. Un cálculo sobre una ventana finita
comprueba las fibras de esa ventana, pero no implica por sí solo el
cuantificador sobre todas las profundidades.
\end{resultbox}

\begin{theorem}[Composición de la dinámica interna]
\label{thm:composicion-dinamica-tpk-rev11}
Una transición observable del TPK es la composición tipada
\[
 \operatorname{Tra}_t(y,d):=\operatorname{Lift}(T_d)(y),
 \qquad
 \operatorname{Upd}_t:=\operatorname{Rec}_t\circ\operatorname{Mem}_t,
\]
\[
 \mathcal U_t(x)
 =\operatorname{Upd}_t\!\left(
  \operatorname{Tra}_t\!\left(
   \operatorname{Sel}_t\!\left(x,M_{\rm ph}(g(t))\right),d(t)
  \right)\right)
\]
sobre \(\mathcal X_{\rm TPK}^{\rm enr}\), seguida, cuando corresponde, por
un lector modular, una extensión cilíndrica o una de las dos proyecciones.
Cada factor actúa sobre un componente distinto del estado; por ello el
selector, el movimiento, el reloj, la memoria y la publicación no pueden
colapsarse en una sola reducción modular.
\end{theorem}

La tabla anterior enumera las familias operatorias. Las equivalencias entre
realizaciones locales deben entrelazar sus dominios, reglas e invariantes.
Los desarrollos posteriores seleccionan composiciones de estas operaciones;
no introducen un cuarto objeto primitivo junto a APP, TRIT y TPK.
